\documentclass[10pt,a4paper]{article}
\usepackage[left=2.5cm,right=2.5cm,top=2.5cm,bottom=2.5cm,]{geometry}
\usepackage{amsmath,amssymb,amsthm}
\newtheorem{theorem}{Theorem}[section]
\newtheorem{lemma}[theorem]{Lemma}
\newtheorem{proposition}[theorem]{Proposition}
\usepackage{setspace}
\setlength{\parindent}{0pt}
\onehalfspacing

\begin{document}
\subsubsection{Binary Operation}
Binary Operation is a rule for combining two elements produce to another element. The binary peration
on $S$ is a mapping of the elements of Cartesian product $S\times S\to S$.
\[f:S\times S\to S\]
For each binary operation, it have to satisfy :
\[\forall a,b\in S, f(a,b)\in S\land f(a,b)\:is\:unique\]
Let's define the binary operation $*$ on $S$, we say $*$ it satisfy the communicative law if and only if :
\[\forall a,b\in S, a*b=b*a\]
we say it satisfy the associative law :
\[\forall a,b,c\in S (a*b)*c=a*(b*c)\]
we say it has neutral $e$ :
\[\forall a\in S, a*e=e*a\]
we say it has inverse :
\[\forall a*a^{-1}=e\]

suppose: $A\cap B=\varnothing$
\[A+e=(A-e)\cup(e-A)=A\cup e=A\]
\[e\subsetneq A\] suppose is wrong\\
otherwise:\\
\[(A-e)\cup(e-A)=A\]
\[(A-A\cap e)\cup(e-A\cap e)=A\]
\[2(A\cap e)=e\]


\end{document}

